**UPDATES:**

**Bought knitting needles and yarn to learn**

**NOTES:**

NAT: "Their basic idea: given a machine knitting program (in knitout), how much yarn will the program use up? This is roughly the same as how "stretched-out" the knit object is -- if you have 50 x 50 stitches, and you use up Y length of yarn, you can sorta-kinda calculate how much yarn is between each stitch."

**NOTES:**

> “Can we predict how much yarn a knitting program will use _just by reading the code_?”

They say yes — by treating knitting as a sum of small building blocks, each using a fixed amount of yarn.

They break knitting into

- knit → makes a loop
- tuck → stacks yarn without pulling through
- miss → skips a needle
- edge-knit → turning at edges
- diagonals → yarn connecting between needles
- x-knit → crossing beds

Every knit object = a bunch of these repeated operations.

ASSUMED:
each knit uses ~constant yarn  
each tuck uses ~constant yarn  

total yarn 

total_yarn =   
  (#knit × yarn_per_knit)  
+ (#tuck × yarn_per_tuck)  
+ ...

called **linear model**.

finding yarn per operation:

1. Knit several sample swatches
2. Each swatch is designed to emphasize one operation
3. Weigh each swatch
4. Count operations in its knitout

So each swatch gives an equation like:
9480·d1 + 118·edge + 9422·knit = total_mass

They stack a bunch of these and solve:
operation counts x unknown yarn per operation = measured yarn used


This works because mahines are very consistent, they are repeating it thousnads of times and local geometry is the same

So even though fabric is messy the _average yarn per operation becomes stable_

After fitting, they got values like

- knit ≈ 0.39 mg
- tuck ≈ 0.41 mg
- miss ≈ 0.14 mg
- x-knit ≈ 0.84 mg

- loops (knit/tuck) use more yarn
- skips (miss) use less
- cross-bed operations use more


Once they know this “cost per operation”

1. Take a new knitout file
2. Count operations
3. Multiply and sum

predicted yarn usage

They measure mass, not length (because scales are easy).

**length = mass / linear_density**

So they can estimate actual yarn length.

stitch size & yarn type can be tested with stich size: bigger stitches → more yarn per knit. Relationship is roughly linear. and results can be scaled using yarn density

You don’t need a physics simulation of yarn.

Lets us
- estimate cost before knitting
- simulate fabric density
- optimize designs
- even map where yarn is concentrated
- etc.


**How we could replicate this experiment?**

Write knitout swatches (each isolates an operation)
Knit them on machine
Measure yarn used (mass or length)
parse knitout → count ops

1. Write knitout swatches (each isolates an operation)
	- Example knitout snippet:
		knit + f0 1  
		knit - f0 1  
		tuck + f1 1  
		miss + f2 1
	- Make ~5–7 swatches:
		1. All knit → measures knit cost
		2. Tuck pattern → measures tuck cost
		3. Miss pattern (skip needles) → measures miss cost
		4. Narrow strip → exaggerates edge
2. Knit them on machine
3. Measure yarn used (mass or length)
	- Weigh yarn cone before knitting
	- Knit swatch
	- Weigh after
	- Difference = yarn mass used
4. Parse knitout → count ops
	- write py script
	- knit_count, tuck_count, miss_count to calculate total_yarn_used
5. Fit the model, yarn per __
	- operation counts x unknown yarn per operation = measured yarn used


